% GENERATED FILE. Edit canonical lessons, then run npm run generate:books.
% canonical-source-hash: fnv1a64:4892d9167b7e7381
% canonical-lessons: JA-C320-tachiba, JA-C320-gimon, JA-C320-shuchou-suru, JA-C320-rikai-suru, JA-C320-tokuni

\chapter{Seeing Both Sides}
\label{ch:ja-seeing-both-sides}
\hlchaptermodality{320}

\begin{chapteropening}
\textbf{By the end of this chapter:} \emph{I can say how a problem looks from another side.}

Complete the last lesson of chapter 320: I can say how a problem looks from another side.
\end{chapteropening}

\section[\texorpdfstring{\ja{たちば} (tachiba)}{たちば (tachiba)}]{\ja{たちば} (tachiba) — a position; a point of view}

\label{lesson:JA-C320-tachiba}



\begin{quote}
Before the new one: say the Japanese for to misunderstand, then the Japanese for in earnest.
\end{quote}

\subsection*{You'll want to know: \ja{たちば}}
\textbf{\ja{たちば}} — \emph{tachiba} — \textquotedblleft{}a position; a point of view\textquotedblright{}.

\begin{quote}
\textbf{\ja{せんせいの} \ja{たちばも} \ja{かんがえた} \ja{ほうが} \ja{いいです。}} — \emph{sensei no tachiba mo kangaeta hō ga ii desu} — \textquotedblleft{}You should think about the teacher's position too.\textquotedblright{}
\end{quote}

\begin{cousinweb}
\textbf{\ja{たちば}} (\emph{tachiba}) is a native Japanese word (wago), inherited rather than borrowed from Chinese: it is built from \emph{tachi}, \textquotedblleft{}standing (tatsu)\textquotedblright{} and \emph{ba}, \textquotedblleft{}place\textquotedblright{}, and its kanji were attached later as kun'yomi (JMdict).
\end{cousinweb}

\begin{grammarlens}[title={What you've built}]
One more word for giving reasons and opinions.
\end{grammarlens}

\subsection*{Your turn}
\begin{itemize}
\raggedright
  \item \emph{Say it:} \emph{tachiba}
  \item \emph{Say it:} \emph{tachiba}, once more
  \item \emph{Say it:} \emph{sensei no tachiba mo kangaeta hō ga ii desu} — You should think about the teacher's position too.
  \item \emph{Recall:} say the Japanese for to misunderstand, then the Japanese for in earnest, then say \emph{tachiba} again
\end{itemize}

\begin{tcolorbox}[breakable,colback=teal!4,colframe=teal!35!black,title={Before you move on}]
What does \ja{たちば} mean? (\textquotedblleft{}A position; a point of view\textquotedblright{}.) Say it once more.
\end{tcolorbox}
\section[\texorpdfstring{\ja{ぎもん} (gimon)}{ぎもん (gimon)}]{\ja{ぎもん} (gimon) — a doubt; a question in one's mind}

\label{lesson:JA-C320-gimon}



\begin{quote}
Before the new one: say the Japanese for in earnest, then the Japanese for a point of view.
\end{quote}

\subsection*{You'll want to know: \ja{ぎもん}}
\textbf{\ja{ぎもん}} — \emph{gimon} — \textquotedblleft{}a doubt; a question in one's mind\textquotedblright{}.

\begin{quote}
\textbf{\ja{ぎもんが} \ja{あるので、しつもんします。}} — \emph{gimon ga aru node, shitsumon shimasu} — \textquotedblleft{}I have a doubt, so I will ask a question.\textquotedblright{}
\end{quote}

\begin{cousinweb}
\textbf{\ja{ぎもん}} (\emph{gimon}) is Sino-Japanese (kango): two Chinese characters read with their on'yomi, readings borrowed from Chinese. They are \emph{gi}, \textquotedblleft{}doubt\textquotedblright{} (Middle Chinese \emph{ngi}) and \emph{mon}, \textquotedblleft{}question\textquotedblright{} (\emph{mjunH}) (KANJIDIC2; Middle Chinese in Baxter's transcription).
\end{cousinweb}

\begin{grammarlens}[title={What you've built}]
One more word for giving reasons and opinions.
\end{grammarlens}

\subsection*{Your turn}
\begin{itemize}
\raggedright
  \item \emph{Say it:} \emph{gimon}
  \item \emph{Say it:} \emph{gimon}, once more
  \item \emph{Say it:} \emph{gimon ga aru node, shitsumon shimasu} — I have a doubt, so I will ask a question.
  \item \emph{Recall:} say the Japanese for in earnest, then the Japanese for a point of view, then say \emph{gimon} again
\end{itemize}

\begin{tcolorbox}[breakable,colback=teal!4,colframe=teal!35!black,title={Before you move on}]
What does \ja{ぎもん} mean? (\textquotedblleft{}A doubt; a question in one's mind\textquotedblright{}.) Say it once more.
\end{tcolorbox}
\section[\texorpdfstring{\ja{しゅちょうする} (shuchō suru)}{しゅちょうする (shuchō suru)}]{\ja{しゅちょうする} (shuchō suru) — to insist; to claim}

\label{lesson:JA-C320-shuchou-suru}



\begin{quote}
Before the new one: say the Japanese for a point of view, then the Japanese for a doubt.
\end{quote}

\subsection*{You'll want to know: \ja{しゅちょうする}}
\textbf{\ja{しゅちょうする}} — \emph{shuchō suru} — \textquotedblleft{}to insist; to claim\textquotedblright{}.

\begin{quote}
\textbf{\ja{ただしいと} \ja{おもったから、つよく} \ja{しゅちょうしました。}} — \emph{tadashii to omotta kara, tsuyoku shuchō shimashita} — \textquotedblleft{}He thought he was right, so he insisted strongly.\textquotedblright{}
\end{quote}

Polite \textbf{\ja{しゅちょうします}}, \emph{shuchōshimasu}; a \ja{する} verb. Past (polite): \textbf{\ja{おとうとが} \ja{しゅちょうしました。}}, \emph{otōto ga shuchō shimashita}.

\begin{cousinweb}
\textbf{\ja{しゅちょうする}} (\emph{shuchō suru}) is built from a Sino-Japanese noun and the native verb \emph{suru}, \textquotedblleft{}to do\textquotedblright{}. The noun, \emph{shuchō}, is two Chinese characters read with their on'yomi, readings borrowed from Chinese: \emph{shu}, \textquotedblleft{}main thing\textquotedblright{} (Middle Chinese \emph{tsyuX}) and \emph{chō}, \textquotedblleft{}stretch\textquotedblright{} (\emph{trjang}) (KANJIDIC2; JMdict; Middle Chinese in Baxter's transcription).
\end{cousinweb}

\begin{grammarlens}[title={What you've built}]
One more word for giving reasons and opinions.
\end{grammarlens}

\subsection*{Your turn}
\begin{itemize}
\raggedright
  \item \emph{Say it:} \emph{shuchō suru}
  \item \emph{Say it:} \emph{shuchō suru}, once more
  \item \emph{Say it:} \emph{tadashii to omotta kara, tsuyoku shuchō shimashita} — He thought he was right, so he insisted strongly.
  \item \emph{Recall:} say the Japanese for a point of view, then the Japanese for a doubt, then say \emph{shuchō suru} again
\end{itemize}

\begin{tcolorbox}[breakable,colback=teal!4,colframe=teal!35!black,title={Before you move on}]
What does \ja{しゅちょうする} mean? (\textquotedblleft{}To insist; to claim\textquotedblright{}.) Say it once more.
\end{tcolorbox}
\section[\texorpdfstring{\ja{りかいする} (rikai suru)}{りかいする (rikai suru)}]{\ja{りかいする} (rikai suru) — to comprehend; to understand fully}

\label{lesson:JA-C320-rikai-suru}



\begin{quote}
Before the new one: say the Japanese for a doubt, then the Japanese for to insist.
\end{quote}

\subsection*{You'll want to know: \ja{りかいする}}
\textbf{\ja{りかいする}} — \emph{rikai suru} — \textquotedblleft{}to comprehend; to understand fully\textquotedblright{}.

\begin{quote}
\textbf{\ja{くわしく} \ja{せつめいしたので、みんな} \ja{りかいしました。}} — \emph{kuwashiku setsumei shita node, minna rikai shimashita} — \textquotedblleft{}I explained it in detail, so everyone understood.\textquotedblright{}
\end{quote}

Polite \textbf{\ja{りかいします}}, \emph{rikaishimasu}; a \ja{する} verb. Past (polite): \textbf{\ja{みんなが} \ja{りかいしました。}}, \emph{minna ga rikai shimashita}.

\begin{cousinweb}
\textbf{\ja{りかいする}} (\emph{rikai suru}) is built from a Sino-Japanese noun and the native verb \emph{suru}, \textquotedblleft{}to do\textquotedblright{}. The noun, \emph{rikai}, is two Chinese characters read with their on'yomi, readings borrowed from Chinese: \emph{ri}, \textquotedblleft{}reason\textquotedblright{} (Middle Chinese \emph{liX}) and \emph{kai}, \textquotedblleft{}understanding\textquotedblright{} (\emph{keaX}) (KANJIDIC2; JMdict; Middle Chinese in Baxter's transcription).
\end{cousinweb}

\begin{grammarlens}[title={What you've built}]
One more word for giving reasons and opinions.
\end{grammarlens}

\subsection*{Your turn}
\begin{itemize}
\raggedright
  \item \emph{Say it:} \emph{rikai suru}
  \item \emph{Say it:} \emph{rikai suru}, once more
  \item \emph{Say it:} \emph{kuwashiku setsumei shita node, minna rikai shimashita} — I explained it in detail, so everyone understood.
  \item \emph{Recall:} say the Japanese for a doubt, then the Japanese for to insist, then say \emph{rikai suru} again
\end{itemize}

\begin{tcolorbox}[breakable,colback=teal!4,colframe=teal!35!black,title={Before you move on}]
What does \ja{りかいする} mean? (\textquotedblleft{}To comprehend; to understand fully\textquotedblright{}.) Say it once more.
\end{tcolorbox}
\section[\texorpdfstring{\ja{とくに} (tokuni)}{とくに (tokuni)}]{\ja{とくに} (tokuni) — especially; in particular}

\label{lesson:JA-C320-tokuni}



\begin{quote}
Before the new one: say the Japanese for to insist, then the Japanese for to comprehend.
\end{quote}

\subsection*{You'll want to know: \ja{とくに}}
\textbf{\ja{とくに}} — \emph{tokuni} — \textquotedblleft{}especially; in particular\textquotedblright{}.

\begin{quote}
\textbf{\ja{なつは} \ja{とくに} \ja{あついから、みずを} \ja{のんで} \ja{ください。}} — \emph{natsu wa tokuni atsui kara, mizu o nonde kudasai} — \textquotedblleft{}Summer is especially hot, so please drink water.\textquotedblright{}
\end{quote}

\begin{cousinweb}
\textbf{\ja{とくに}} (\emph{tokuni}) is Sino-Japanese (kango): one Chinese character read with its on'yomi, reading borrowed from Chinese. It is \emph{toku}, \textquotedblleft{}special\textquotedblright{} (Middle Chinese \emph{dok}) (KANJIDIC2; Middle Chinese in Baxter's transcription). The native particle \emph{ni} after it makes the adverb (JMdict).
\end{cousinweb}

\begin{grammarlens}[title={What you've built}]
One more word for giving reasons and opinions.
\end{grammarlens}

\subsection*{Your turn}
\begin{itemize}
\raggedright
  \item \emph{Say it:} \emph{tokuni}
  \item \emph{Say it:} \emph{tokuni}, once more
  \item \emph{Say it:} \emph{natsu wa tokuni atsui kara, mizu o nonde kudasai} — Summer is especially hot, so please drink water.
  \item \emph{Recall:} say the Japanese for to insist, then the Japanese for to comprehend, then say \emph{tokuni} again
\end{itemize}

\begin{tcolorbox}[breakable,colback=teal!4,colframe=teal!35!black,title={Before you move on}]
What does \ja{とくに} mean? (\textquotedblleft{}Especially; in particular\textquotedblright{}.) Say it once more.
\end{tcolorbox}
